% 首测可行集、推荐点与候选凸包的统一示意图。
% 左侧给出全局位置，右侧对左侧圈出的候选凸包作局部放大。
\begin{tikzpicture}[x=1cm,y=1cm,>=Stealth,font=\small]
  \path[use as bounding box] (0,0) rectangle (13.6,5.48);

  \node at (3.0,.20) {(a) 首测可行集与推荐区域};
  \node at (10.55,.20) {(b) 候选凸包的局部放大};

  % -------------------- 左侧全局图 --------------------
  \begin{scope}[shift={(3.0,2.75)},x=.62cm,y=.62cm]
    \draw[gray!70,dashed] (0,0) circle (3.600);
    \draw[gray!55,dotted] (0,0) circle (3.000);
    \draw[->,black] (-2.2,0)--(3.4,0)
      node[right,font=\small] {$x$};
    \draw[->,black] (0,-1.7)--(0,1.3)
      node[above,font=\small] {$y$};

    \node[gray!85,font=\scriptsize,anchor=south west]
      at (-2.9,-2.9) {目标圆域半径 $1800\,\mathrm{m}$};

    % 首测可行源集采用浅蓝填充、深蓝边界。
    \path[fill=qCyanLight,draw=qBlueDark,line width=1.1pt]
      (0,0)--(3,-0.05236)--(3,-0.04189)--(3,-0.02094)--
      (3,0)--(3,0.02094)--(3,0.04189)--(3,0.05236)--cycle;
    \node[text=qBlueDark,font=\small,anchor=west] at (2.15,0.48) {$K_1$};

    \fill (0,0) circle (0.035);
    \node[below left,font=\small] at (0,0) {$S_1$};

    % 候选凸包在全局图中的位置。
    \path[fill=qBlueLight,draw=qBlueDark,line width=.8pt]
      (1.519,-1.291)--(1.529,-1.291)--(1.569,-1.251)--
      (1.629,-1.171)--(1.679,-1.101)--(1.729,-1.021)--
      (1.779,-0.9312)--(1.819,-0.8512)--(1.819,-0.8412)--
      (1.789,-0.8712)--(1.739,-0.9312)--(1.619,-1.091)--
      (1.559,-1.181)--(1.519,-1.251)--cycle;
    \fill[qBlueDark] (1.679,-1.101) circle (0.045);
    \node[text=black,font=\small,anchor=north west]
      at (1.76,-1.13) {$\widehat q$};
    \draw[black,densely dashed] (0,0)--(1.679,-1.101);
    \node[text=black,font=\scriptsize,anchor=north]
      at (-0.45,-1.55) {$1004\,\mathrm{m}$};
  \end{scope}

  % 左侧放大标记及通往右侧面板的虚线。
  \draw[qBlueDark,line width=.8pt] (4.041,2.067) circle (.24);
  \draw[gray!70,dashed] (4.21,2.24)--(8.30,4.49);
  \draw[gray!70,dashed] (4.21,1.90)--(8.30,1.00);

  % -------------------- 右侧局部放大图 --------------------
  % 横坐标换算：X=(x-735)*0.0225 cm；纵坐标换算：Y=(y+675)*0.0125 cm。
  \begin{scope}[shift={(8.30,1.00)},x=.0225cm,y=.0125cm]
    \foreach \xx/\lab in {0/735,50/785,100/835,150/885,200/935}{
      \draw[gray!20,thin] (\xx,0)--(\xx,279);
      \node[font=\scriptsize,gray!85,anchor=north] at (\xx,-8) {\lab};
    }
    \foreach \yy/\lab in
      {25/-650,75/-600,125/-550,175/-500,225/-450,275/-400}{
      \draw[gray!20,thin] (0,\yy)--(200,\yy);
      \node[font=\scriptsize,gray!85,anchor=east] at (-7,\yy) {\lab};
    }

    \draw[gray!55,thin] (0,0) rectangle (200,279);

    \path[fill=qBlueLight,draw=qBlueDark,line width=.9pt]
      (24.707,29.395)--(29.707,29.395)--(49.707,49.395)--
      (79.707,89.395)--(104.707,124.395)--(129.707,164.395)--
      (154.707,209.395)--(174.707,249.395)--(174.707,254.395)--
      (159.707,239.395)--(134.707,209.395)--(74.707,129.395)--
      (44.707,84.395)--(24.707,49.395)--cycle;

    \fill[qBlueDark] (104.707,124.395) circle (1.8);
    \node[text=black,font=\small,anchor=west]
      at (114.707,124.395) {$\widehat q$};

    \node[text=black,font=\scriptsize,anchor=north west,align=left]
      at (14,263) {候选凸包：$14$ 个顶点\\等效宽度约 $20\,\mathrm{m}$};
    \node[font=\scriptsize,gray!85,anchor=north west] at (210,-7.8) {$x/\mathrm{m}$};
    \node[font=\scriptsize,gray!85,anchor=south east] at (-8,-20) {$y/\mathrm{m}$};
  \end{scope}
\end{tikzpicture}
